\booksection{الفصل 20: وضع المطور والصلاحيات المتقدمة}

\chapterintro{هذا الفصل مُعد للمستخدمين الذين يرغبون في التعمق في “الوظائف المخفية” للهاتف. وضع المطور، وصلاحية Root، وقفل BL... هذه المفاهيم التي تبدو غريبة، ليست بعيدة عن المستخدم العادي كما يظن.}

\booksubsection{أولاً: وضع المطور: افتح الإعدادات المخفية}

وضع المطور قائمة مخفية مدمجة في نظام الهاتف، تحتوي على وظائف لا يحتاجها المستخدم العادي لكن المطورين يحتاجونها.

\subsubsection*{كيفية التفعيل:}

\begin{itemize}
  \item \textbf{أندرويد}: ادخل الإعدادات > حول الهاتف > انقر "رقم الإصدار" 7 مرات متتالية، حتى تظهر رسالة "أنت الآن في وضع المطور".
  \item \textbf{HarmonyOS}: الإعدادات > حول الهاتف > انقر "إصدار HarmonyOS" 7 مرات متتالية.
  \item \textbf{iOS}: لا يمكن للمستخدم العادي تفعيل وضع المطور الكامل، يتطلب حساب Apple Developer مدفوع.
\end{itemize}

\subsubsection*{نظرة على الوظائف الشائعة:}

\begin{itemize}
  \item \textbf{تصحيح USB}: مفتاح أساسي لربط الحاسوب لإجراء عمليات متقدمة (مثل الوميض واستعادة البيانات وأوامر ADB)، انظر قسم "أوامر ADB العملية" في هذا الفصل.
  \item \textbf{إظهار عمليات اللمس}: عرض نقاط اللمس على الشاشة (نقطة صغيرة تتبع الإصبع)، مناسب لتسجيل شروحات العمليات أو العروض.
\end{itemize}

\subsubsection*{متعلق بالأداء والسلاسة:}

\begin{itemize}
  \item \textbf{ضبط سرعة الرسوم المتحركة (مقياس رسوم النوافذ / مقياس رسوم الانتقال / مقياس مدة برنامج الرسوم)}: تتحكم هذه الخيارات الثلاثة في مدة الرسوم المتحركة عند فتح التطبيقات وتبديل الصفحات. الافتراضي 1x، وضبطه على 0.5x يجعل الهاتف "يبدو" أسرع (البرنامج لا يعمل أسرع فعلاً، فقط الرسوم أسرع)، والإيقاف يقفز مباشرة بلا رسوم. معظم الخبراء يضبطونه على 0.5x للحفاظ على الإحساس بالانتقال مع ظهور السرعة؛ الإيقاف الكامل يجعل العملية تبدو جامدة وقد يخفي التلعثم الحقيقي للبرنامج، مما يجعلك تعتقد خطأً أنه "سلس جداً".
  \item \textbf{إجبار عرض GPU (Force GPU Rendering)}: يجبر جميع رسومات 2D على المعالجة عبر GPU بدلاً من CPU. في الموديلات القديمة أو بعض التطبيقات غير المُحسّنة يمكن تحسين سلاسة التمرير، لكنه يزيد حمل GPU واستهلاك الطاقة. الهواتف الحديثة تستخدم GPU بشكل كبير افتراضياً، وهذا المفتاح محدود التأثير على الموديلات المتوسطة والعالية، وهو خيار تاريخي في الغالب، المستخدم العادي يبقي على الافتراضي.
  \item \textbf{عرض تحليل وضع عرض GPU (Profile GPU Rendering)}: يعرض على الشاشة مخططاً بالأعمدة لزمن عرض كل إطار في الوقت الفعلي، الخط الأخضر يمثل 16.6ms (عتبة زمن الإطار لـ 60fps). الأعمدة التي تتجاوز الخط الأخضر تعني إسقاط إطارات. هذه أداة للمطورين لتشخيص تلعثم التطبيقات، ويمكن للمستخدم العادي استخدامها لمقارنة أي تطبيق سيئ التحسين — مثلاً تطبيق تتجاوز أعمدته الخط الأحمر كثيراً عند التمرير، فهو المسؤول غالباً.
\end{itemize}

\subsubsection*{متعلق بالخلفية والبطارية:}

\begin{itemize}
  \item \textbf{تقييد العمليات الخلفية}: يمكن ضبط "الحد القياسي" أو "عملية واحدة كحد أقصى" أو "عمليتين كحد أقصى" وغيرها. يبدو أن تقييد العمليات الخلفية يوفر الطاقة، لكن إدارة الخلفية في Android جيدة بالفعل، والتقييد القسري يجعل التطبيقات تُعاد تحميلها بشكل متكرر (التشغيل البارد يستهلك طاقة أكثر وتجربة أسوأ). هذا الخيار أنسب للتصحيح، ولا يُنصح بفتحه بشكل دائم للمستخدم العادي، اتركه على "الحد القياسي".
  \item \textbf{عدم الاحتفاظ بالأنشطة}: يُدمّر كل نشاط بمجرد مغادرة الواجهة. غير موصى به على الإطلاق، يؤدي لبياض الشاشة المتكرر وإعادة تحميل التطبيقات، يُستخدم فقط لاختبار منطق استعادة التطبيقات.
\end{itemize}

\begin{tcbwarn}{تحذيرات}
إعدادات وضع المطور قد تؤثر على استقرار النظام، لا تغيّر ما لا تفهمه. بعد تفعيل تصحيح USB، كن حذراً عند الاتصال بالحاسوب، لا تتصل بحواسيب غير موثوقة.

\end{tcbwarn}

\booksubsection{اثنان: OEM Unlock: “مفتاح الوميض” في خيارات المطور}

في خيارات المطور هناك مفتاح غير لافت لكن مهم — “OEM Unlock” (تسميه بعض العلامات “السماح بفتح OEM” أو “فتح قفل الجهاز”). هذا المفتاح نفسه لا يفتح قفل BL، لكنه “ترخيص مسبق”: فقط بفتحه تسمح الشركة بفتح BL عبر أوامر fastboot أو أدوات العلامة.

\begin{itemize}
  \item \textbf{وظيفة المفتاح}: يعادل إعلان الشركة "أوافق على التنازل عن الضمان وتحمل المخاطر، اسمح لي بفتح BL". إذا كان هذا المفتاح مغلقاً، تُرفض جميع أوامر فتح BL مباشرة.
  \item \textbf{العلاقة مع قفل BL}: OEM Unlock هو "تأهيل الطلب"، وفتح BL هو "العملية الفعلية". الأول نقرة في إعدادات الهاتف، والثاني يتطلب طلب رمز فتح من موقع الشركة وتشغيل أوامر fastboot على الحاسوب (مثل\texttt{fastboot flashing unlock})، وكلاهما لا غنى عنه.
  \item \textbf{هل يعيد الاستعادة للمصنع ضبطه؟}: نعم. يعود مفتاح "OEM Unlock" إلى الحالة المغلقة بعد الاستعادة للمصنع، لكن قفل BL بمجرد فتحه يكون حالة على مستوى العتاد ولا يُغلق بالاستعادة للمصنع — يجب تنفيذ\texttt{fastboot flashing lock}يدوياً لإعادة القفل.
  \item \textbf{متى يجب فتحه}: إلا إذا كنت متأكداً من الوميض أو Root، أبقه مغلقاً. هو و"تصحيح USB" مفتاحان مستقلان، وعادةً يحتاج فتحهما معاً عند الوميض.
\end{itemize}

\begin{tcbinfo}{معلومة: لماذا تضع الشركات هذا المفتاح}
في Android المبكر لم يكن هذا المفتاح موجوداً، والبرمجيات الخبيثة بمجرد حصولها على الصلاحية يمكنها وميض النظام سراً.لاحقاً طالبت Google جميع شركات Android بإضافة مفتاح “OEM Unlock” في خيارات المطور، كتأكيد “موافقة المستخدم المعلنة”. لهذا حتى لو فتحت قفل BL، إذا كان مفتاح OEM Unlock مغلقاً، سيفشل الوميض مرة أخرى.

\end{tcbinfo}

\booksubsection{ثلاثة: قفل BL (Bootloader): الباب الأول للوميض}

Bootloader هو البرنامج الذي يحمّل نظام التشغيل عند إقلاع الهاتف. قفل BL هو قفل تضعه الشركات المصنّعة للأمان، لمنع المستخدمين من تفليش أنظمة غير رسمية.

\subsubsection*{وظائف قفل BL:}

\begin{itemize}
  \item \textbf{حماية أمنية}: منع البرمجيات الخبيثة من تفليش النظام.
  \item \textbf{مضاد للسرقة}: حتى لو فُلش الهاتف، لا يمكن استخدامه بدون كلمة مرور الحساب الأصلية.
  \item \textbf{تقييد الضمان}: فتح قفل BL عادةً يعني فقدان الضمان الرسمي.
\end{itemize}

\subsubsection*{سياسات قفل BL لدى العلامات التجارية:}

\begin{itemize}
  \item \textbf{Xiaomi}: يمكن فتحه، يتطلب ربط الحساب وفترة انتظار (7-30 يوماً).
  \item \textbf{Huawei}: توقفت عن تقديم رمز فتح القفل بعد 2018، يصعب فتحه رسمياً.
  \item \textbf{OnePlus}: يمكن فتحه، سياسة متساهلة نسبياً.
  \item \textbf{OPPO/vivo}: بعض الموديلات يمكن فتحها، تتطلب تقديم طلب.
  \item \textbf{Samsung}: يمكن فتحه، لكنه يُفعّل فصل Knox، فتفقد Samsung Pay وغيرها من الوظائف.
  \item \textbf{Apple}: أجهزة iOS لا تحتوي على مفهوم قفل BL بالمعنى التقليدي، لكن لها آليات أمان مشابهة.
\end{itemize}

\booksubsection{أربعة: أوامر ADB العملية: التحكم في الهاتف عبر الحاسوب}

ADB اختصار لـ Android Debug Bridge (جسر تصحيح Android)، أداة سطر أوامر رسمية من Google، تتيح للحاسوب التواصل مع هاتف Android عبر USB أو Wi-Fi. لا تحتاج Root نفسه، يكفي فتح “تصحيح USB” على الهاتف. يمكن فهمه كـ “جهاز تحكم عن بعد تصدر أوامر للهاتف عبر الحاسوب” — الوميض وتثبيت التطبيقات ونقل الملفات وقراءة السجلات تعتمد عليه.

\subsubsection*{كيفية الاتصال عبر USB:}

\begin{enumerate}
  \item جانب الهاتف: ادخل خيارات المطور افتح "تصحيح USB".
  \item جانب الحاسوب: نزّل Platform Tools من موقع Google الرسمي (يحتوي\texttt{adb}و\texttt{fastboot}ملفين قابلين للتنفيذ)، وفك الضغط في أي مجلد.
  \item اربط الهاتف بالحاسوب بكابل بيانات، سيظهر تنبيه على الهاتف "السماح بتصحيح USB؟"، اختر "السماح دائماً" ثم موافق.
  \item في سطر أوامر الحاسوب (PowerShell أو cmd في Windows) نفّذ\texttt{adb devices}، إذا كان المخرج مشابهاً لـ\texttt{List of devices attached /\allowbreak  xxxxxxxx device}فهذا يعني نجاح الاتصال.
\end{enumerate}

\begin{tcbwarn}{لا تتجاهل التعريف والتنبيه}
إذا عرض مستخدم Windows\texttt{adb devices}“unauthorized” أو لم يدرج الجهاز أصلاً، فالسبب غالباً عدم تثبيت التعريف أو نسيان الضغط على “السماح” على الهاتف. Samsung/Xiaomi وغيرها يمكن تثبيت “USB Driver” أو “مساعد الهاتف” الخاص بها لتكملة التعريف؛ Mac و Linux عموماً بلا تعريف. إذا أغلقت التنبيه بالخطأ، افصل وأعد كابل البيانات ليظهر مرة أخرى.

\end{tcbwarn}

\subsubsection*{أوامر ADB الشائعة:}

\begin{itemize}
  \item \texttt{adb install app.\allowbreak apk}: تثبيت APK من الحاسوب على الهاتف (\texttt{-\allowbreak r}يعني التثبيت فوق النسخة مع الاحتفاظ بالبيانات؛\texttt{-\allowbreak s}تثبيت على SD card).
  \item \texttt{adb uninstall com.\allowbreak xxx.\allowbreak xxx}: إزالة تثبيت تطبيق باسم الحزمة.
  \item \texttt{adb push مسار\_الحاسوب /\allowbreak sdcard/\allowbreak المسار\_المستهدف}: دفع ملف من الحاسوب إلى الهاتف.
  \item \texttt{adb pull /\allowbreak sdcard/\allowbreak ملف مسار\_الحاسوب}: سحب ملف من الهاتف إلى الحاسوب (أسرع بكثير من الإرسال عبر WeChat ثم التنزيل).
  \item \texttt{adb shell}: دخول Linux Shell في الهاتف، يمكن تنفيذ أوامر\texttt{ls}و\texttt{pm}و\texttt{am}مباشرة.
  \item \texttt{adb reboot bootloader}: إعادة تشغيل الهاتف في وضع Fastboot (مدخل الوميض).
  \item \texttt{adb reboot recovery}: إعادة تشغيل في وضع Recovery.
  \item \texttt{adb shell pm list packages}: سرد أسماء حزم جميع التطبيقات المثبتة على الهاتف (ضروري لإيجاد اسم الحزمة). أضف\texttt{-\allowbreak s}لعرض تطبيقات النظام فقط، أضف\texttt{-\allowbreak 3}لعرض تطبيقات الطرف الثالث فقط.
  \item \texttt{adb shell pm uninstall -\allowbreak -\allowbreak user 0 اسم\_الحزمة}: إزالة تثبيت تطبيق مُثبّت مسبقاً للمستخدم الحالي، بدون Root، انظر التطبيق العملي أدناه.
\end{itemize}

\subsubsection*{تطبيق عملي: إزالة التطبيقات المُثبّتة مسبقاً بدون Root:}

هذه واحدة من أكثر وظائف ADB فائدة: باستخدام\texttt{pm uninstall -\allowbreak -\allowbreak user 0}يمكن إزالة تثبيت التطبيقات المُثبّتة مسبقاً لـ “المستخدم الحالي” دون الحاجة لـ Root ودون التأثير على قسم النظام. التطبيقات المُزالة “تختفي” من النظام، لكنها لا تزال في قسم النظام فعلياً، وتعود بعد الاستعادة للمصنع.

\begin{enumerate}
  \item أولاً استخدم\texttt{adb shell pm list packages -\allowbreak s}لسرد جميع تطبيقات النظام، وابحث عن اسم حزمة التطبيق الذي تريد إزالته. مثلاً متصفح Xiaomi هو\texttt{com.\allowbreak android.\allowbreak browser}، ومتجر تطبيقات Huawei هو\texttt{com.\allowbreak vmall.\allowbreak client}.
  \item غير متأكد من التطبيق المقابل لاسم الحزمة؟ يمكن استخدام\texttt{adb shell pm list packages | findstr كلمة\_مفتاحية}للبحث على الحاسوب، أو تثبيت تطبيق "Package Viewer" لعرض اسم حزمة كل تطبيق مباشرة على الهاتف.
  \item تنفيذ الإزالة:\texttt{adb shell pm uninstall -\allowbreak -\allowbreak user 0 اسم\_الحزمة}، مثلاً\texttt{adb shell pm uninstall -\allowbreak -\allowbreak user 0 com.\allowbreak android.\allowbreak browser}. إذا عاد\texttt{Success}فهذا يعني النجاح.
  \item تريد الاستعادة؟ طريقتان: (1) الاستعادة للمصنع (الأكثر شمولاً لكن مزعج)؛ (2) إعادة التثبيت:\texttt{adb shell cmd package install-\allowbreak existing اسم\_الحزمة}.
\end{enumerate}

\textbf{لا تحذف حزم النظام الأساسية بشكل عشوائي}\par

رغم أن “pm uninstall --user 0” طريقة آمنة بدون Root، إلا أن حذف تطبيقات نظام أساسية (مثل\texttt{com.\allowbreak android.\allowbreak systemui}شريط الحالة،\texttt{com.\allowbreak android.\allowbreak phone}الهاتف،\texttt{com.\allowbreak android.\allowbreak settings}الإعدادات) يؤدي لمشاكل خطيرة كانهيار النظام وشاشة سوداء وعدم القدرة على الإقلاع. قبل الحذف ابحث في Google “هل يمكن حذف اسم هذه الحزمة”، والأفضل التقليص عن الخطأ. يُنصح بحذف واحد أو اثنين في كل مرة، والتأكد من عدم وجود مشاكل قبل المتابعة.

\begin{tcbinfo}{معلومة: ماذا يعني --user 0}
بعد Android 4.2 دعم تعدد المستخدمين، لكل مستخدم ID مستقل (المالك 0، المستخدمون الإضافيون من 10).\texttt{-\allowbreak -\allowbreak user 0}يعني “إزالة التثبيت لمالك الجهاز الحالي فقط”، ملف APK في قسم النظام لا يزال موجوداً، فقط لا يُعرض للمالك. لهذا لا يحتاج Root — لم يمس ملفات النظام، فقط عدّل “تكوين المستخدم”. وهذا أكبر فرق عن الإزالة بعد Root: إزالة Root تحذف فعلياً من النظام ولن يحصل عليها أي مستخدم؛ إزالة ADB “تخفي للمستخدم الحالي” فقط.

\end{tcbinfo}

\booksubsection{خمسة: صلاحية Root: “المدير الفائق” في Android}

صلاحية Root تعادل المستخدم الفائق في Linux (تحكم كامل في دليل /، يمكن قراءة وكتابة أي ملف)، وهي أعلى صلاحية في نظام Android. تختلف جوهرياً عن Administrator في Windows: Windows لاعتبارات أمنية يفرض قيوداً كثيرة على المدير (أعلى صلاحية فعلية هي TrustedInstaller)، حتى Administrator لا يستطيع تنفيذ\texttt{sudo rm -\allowbreak rf /\allowbreak *}مباشرة مثل Linux root. بعد الحصول على صلاحية Root، يمكنك تعديل ملفات النظام وإزالة تثبيت التطبيقات المُثبّتة مسبقاً واستخدام التطبيقات المتقدمة التي تتطلب Root.

\subsubsection*{فوائد Root:}

\begin{itemize}
  \item إلغاء تثبيت التطبيقات المثبتة مسبقاً والتي لا يمكن إزالتها.
  \item استخدام تطبيقات الوظائف المتقدمة (مثل Titanium Backup وGreenify).
  \item تعديل واجهة النظام ووظائفه.
  \item تثبيت Recovery طرف ثالث وROM مخصص.
\end{itemize}

\subsubsection*{مخاطر Root:}

\begin{itemize}
  \item \textbf{: مثل مجتمع شاومي ومجتمع OnePlus، بها دروس رسمية وشعبية.}: معظم الشركات لا تضمن الهواتف التي تم عمل Root لها.
  \item \textbf{مخاطر أمنية}: التطبيقات الخبيثة قد تحصل على صلاحية Root، مما يسبب عواقب وخيمة.
  \item \textbf{عدم استقرار النظام}: الأخطاء في العمليات قد تؤدي إلى انهيار النظام.
  \item \textbf{تطبيقات البنوك لا تعمل}: كثير من تطبيقات البنوك والدفع ترفض العمل عند اكتشاف Root.
  \item \textbf{فشل تحديثات OTA}: قد لا تستقبل تحديثات النظام الرسمية بشكل طبيعي.
\end{itemize}

\textbf{تحذير شديد}\par

عملية Root تنطوي على مخاطر، قد تجعل الهاتف قالباً (لا يقلع). إذا كنت مبتدئاً، يُنصح بفهم الموضوع جيداً أولاً، والتأكد من وجود نسخة احتياطية كاملة. اقرأ الشرح بعناية قبل العمل، وتأكد من طراز جهازك وإصدار النظام.

\booksubsection{ستة: SafetyNet وPlay Integrity: لماذا لا تعمل تطبيقات البنوك بعد Root}

كثير من مستخدمي Root واجهوا لحظة انهيار: بعد Root بحماس تثبيت مجموعة وحدات، ثم فتح تطبيق البنك — “جهازك غير آمن، لا يمكن الاستخدام”، وGoogle Pay لا يعمل، وحتى Netflix يرفض التشغيل بدقة عالية. وراء ذلك آليتا مصادقة سلامة الجهاز من Google: SafetyNet القديم وPlay Integrity API الجديد.

\subsubsection*{ما هو SafetyNet / Play Integrity:}

\begin{itemize}
  \item \textbf{SafetyNet Attestation (الإصدار القديم، منذ 2017)}: واجهة برمجة تطبيقات لفحص سلامة الجهاز أطلقتها Google مبكراً، يستخدم التطبيقها للاستعلام من خوادم Google "هل جهاز هذا اجتاز اختبار التوافق CTS، وهل هو في حالة رسمية غير معدّلة". النتائج تنقسم إلى Basic Attestation (تحقق أساسي) و Strong Attestation (تحقق قوي، يعتمد على مفتاح العتاد). أعلنت Google أن SafetyNet سيتوقف تدريجياً بعد 2024.
  \item \textbf{Play Integrity API (الإصدار الجديد، منذ 2022)}: البديل الذي أطلقته Google، أكثر صرامة وشمولاً من SafetyNet. إلى جانب سلامة الجهاز، يشمل "سلامة التطبيق" (هل التطبيق تم العبث به) و"سلامة الحساب/التفويض" (هل حساب المستخدم موثوق من Google). التطبيقات الجديدة يجب أن تنتقل إلى Play Integrity.
  \item \textbf{من يستخدمه}: تطبيقات البنوك والمالية والمدفوعات (مثل Google Pay وبنوك الإنترنت وAlipay/WeChat Pay في بعض السيناريوهات) وخدمات البث (حماية حقوق النقل بدقة عالية لـ Netflix) وبعض ألعاب مكافحة الغش (مثل بعض ألعاب Tencent/NetEase) وتطبيقات المكتب المؤسسية (فحص توافق أجهزة MDM).
\end{itemize}

\subsubsection*{نتائج المصادقة الثلاثية لـ Play Integrity:}

\begin{itemize}
  \item \textbf{MEETS\_DEVICE\_INTEGRITY (سلامة الجهاز ناجحة)}: الجهاز يشغل نظام Android الرسمي، Bootloader غير مفتوح القفل، بدون Root، اجتاز اختبار التوافق CTS. يمكن استخدام تطبيقات البنوك والمدفوعات بشكل طبيعي.
  \item \textbf{MEETS\_BASIC\_INTEGRITY (السلامة الأساسية ناجحة)}: الجهاز موثوق أساساً، لكن قد تكون هناك anomalies طفيفة. بعض التطبيقات تخفض المستوى (مثل Netflix يبقى قابلاً للتشغيل لكن يحد الدقة)، وبعض تطبيقات البنوك الصارمة ترفض مباشرة.
  \item \textbf{MEETS\_STRONG\_INTEGRITY (السلامة القوية ناجحة)}: المستوى الأعلى، يتطلب دعم مفتاح العتاد (مثل شريحة Titan M في سلسلة Google Pixel)، معظم موديلات العلامات التجارية الثالثة حتى بدون Root لا تصل لهذا المستوى.
  \item \textbf{بدون أي علامة (NO\_INTEGRITY)}: الجهاز عليه Root أو مفتوح BL أو عليه ROM خارجي أو مُحدد كمحاكي، أغلب تطبيقات المدفوعات/البنوك ترفض التشغيل مباشرة.
\end{itemize}

\subsubsection*{لماذا تتوقف هذه التطبيقات بعد Root:}

يطلب التطبيق عبر Play Integrity API مصادقة من خوادم Google، تفحص Google عدة مؤشرات للجهاز — حالة قفل Bootloader، تعديل قسم النظام، وجود ملف su الثنائي، وضع SELinux المتساهل وغيرها. أي anomaly يؤدي لخفض مستوى المصادقة المُعاد أو فراغه، وتطبيق البنك يرى “الجهاز غير آمن” فيرفض الخدمة.

\subsubsection*{حلول الإخفاء الحديثة: Magisk Hide / Shamiko / Zygisk:}

للاستمتاع بفوائد Root مع استخدام تطبيقات البنوك، طوّر المجتمع سلسلة حلول “إخفاء حالة Root”. الفكرة الأساسية: جعل برنامج الكشف يرى نظاماً “يبدو طبيعياً” بدلاً من نظام معدّل.

\begin{itemize}
  \item \textbf{فلاش Recovery خارجي}: حل Systemless Root السائد حالياً — لا يعدّل قسم النظام، بل يحقق Root عبر التركيب. ميزة هذا التصميم: احتمال فشل تحديث OTA أقل، والإزالة أنظف، وهو أساس جميع حلول الإخفاء اللاحقة. لكن استخدام Magisk لا يزال يتطلب فتح قفل BL.
  \item \textbf{Magisk Hide (توقف التحديث)}: وحدة الإخفاء المدمجة في Magisk مبكراً، تخفي Root عن تطبيقات محددة. لكن منذ 2021 أزالها Magisk رسمياً، بسبب تطور Google المستمر لطرق الكشف، ودخول الطرفين في "لعبة القط والفأر".
  \item \textbf{Zygisk}: آلية جديدة أدخلها Magisk 24.0+، مشابهة لـ Riru، تحقق تحميل الوحدات عبر حقن عملية Zygote، أكثر خفاءً من تعديل النظام مباشرة. تفعيل Zygisk شرط مسبق لاستخدام حلول الإخفاء الحديثة.
  \item \textbf{Shamiko}: وحدة الإخفاء السائدة حالياً، تعمل مع Zygisk، تخفي Root عن تطبيقات محددة (وضع "القائمة البيضاء العكسية") وتخفي Magisk نفسه وتزوّر حالة قفل Bootloader. يُوزّع Shamiko عبر مستودع LSPosed أو GitHub.
  \item \textbf{Play Integrity Fix(osm0sis)}: وحدة Magisk مخصصة للكشف عن Google Play Integrity، تعدّل بصمة الجهاز (fingerprint) ليعتقد Google أن جهازك موديل رسمي بدون Root.
\end{itemize}

\begin{tcbwarn}{لعبة القط والفأر لن تنتهي أبداً}
Google تحدّث قواعد الكشف بشكل دوري، وشركات التطبيقات تتابع بترقية تحقق Play Integrity. حلول الإخفاء التي تعمل اليوم قد تُكتشف غداً. إذا كنت تعتمد كلياً على تطبيقات البنوك/المدفوعات (مثل التحويلات اليومية والدفع بالمسح)، فالخيار الأكثر أماناً: لا تفتح Root، لا تفتح قفل BL. حلول الإخفاء أنسب للاعبين “الهاتف الرئيسي بدون Root + الهاتف الاحتياطي للتجربة”.

\end{tcbwarn}

\booksubsection{سبعة: نصائح للمستخدم العادي}

\begin{itemize}
  \item \textbf{وضع المطور}: يمكن تفعيله، لتجربة تسريع الرسوم وتحليل عرض GPU وغيرها، لكن لا تغيّر خيارات أخرى لا تفهمها.
  \item \textbf{OEM Unlock وقفل BL}: إذا كنت لا تنوي الوميض أو Root، أبقِ المفتاحين على الافتراضي (مغلق) للأمان.
  \item \textbf{ADB}: أداة قوية بدون Root، يستطيع المستخدم العادي استخدامها — مثلاً بـ\texttt{pm uninstall -\allowbreak -\allowbreak user 0}لإزالة التطبيقات المُثبّتة مسبقاً. لكن تأكد قبل الحذف من إمكانية حذف اسم الحزمة.
  \item \textbf{ذكاء اصطناعي يعمل محليًا على الهاتف}: لا يُنصح المستخدم العادي بـ Root، المخاطر أكبر من الفوائد. إذا كنت تحتاجبعض وظائف، يمكنك أولاً معرفة إن كان هناك بديل لا يحتاج Root (مثل ADB لإزالة المُثبّت مسبقاً، أو Shizuku بديل لبعض وظائف Root).
\end{itemize}

\begin{tcbpractice}{جرّب: فعّل وضع المطور واضبط سرعة الرسوم المتحركة}
\begin{enumerate}
  \item ادخل الإعدادات > حول الهاتف.
  \item انقر "رقم الإصدار" 7 مرات متتالية، حتى تظهر رسالة تفعيل وضع المطور.
  \item عُد إلى الإعدادات، وابحث عن "خيارات المطور" (عادةً في أسفل الإعدادات أو ضمن "النظام").
  \item ابحث عن "مقياس حركة النافذة"، "مقياس حركة الانتقال"، "مقياس مدة حركة البرنامج"، واضبطها جميعاً على 0.5x أو أوقفها.
  \item شعر إن كان هاتفك أصبح أسرع.
\end{enumerate}

\end{tcbpractice}
